%!TEX root = tese.tex
%%%% Capa
\begin{titlepage}
  \thispagestyle{empty}
\noindent {\rule[-1ex]{16cm}{0.05cm}}
\begin{center}
\begin{minipage}[s]{1.5cm}
\includegraphics[width=0.6in]{figuras/unicamp.pdf}
\end{minipage}\begin{minipage}[s]{11cm}\noindent
{\begin{center} {\sffamily \Large Universidade Estadual de Campinas }\\ 
{ \sffamily Instituto de
Matem\'atica, Estat\'istica e Computa\c c\~ao Cient\'ifica}\protect \\
{\sffamily Programa de Pós-graduação em Matem\'atica Aplicada}
\end{center}}
\end{minipage}
\begin{minipage}[s]{1.5cm} 
\includegraphics[width=0.6in]{figuras/logo.pdf}
\end{minipage} 
\end{center}
\noindent{\rule[.2ex]{16cm}{0.03cm}}

\vspace{2cm} 
%\font\fontGrande=cmcsc10 scaled 2500
%\font\pessoal=cmr9 scaled 2500 

\begin{center}
\noindent{\LARGE {\sffamily \bfseries 2o Exame de Qualificação:  \\ Escolha
adiada de parâmetros em métodos de pontos \\[2mm] interiores para programação linear}}
\end{center}

\begin{center} 
  
\normalsize \vspace{10mm}

\setcounter{footnote}{1} \vspace{15mm} 
\renewcommand{\thefootnote}{\fnsymbol{footnote}}
{\Large {\bfseries \sffamily Luiz Rafael dos Santos}}

{\sffamily \normalsize{Doutorado em Matem\'atica Aplicada}}
\end{center}
\vspace{1cm}
%\setcounter{footnote}{1}
%\renewcommand{\thefootnote}{\fnsymbol{footnote}}
\begin{center}
 
{ \sffamily Orientador: {\large Prof. Dr.  Aurelio Ribeiro Leite de Oliveira}}\\[2mm] 
%Usar se tiver co-orientador
{ \sffamily Co-orientadores: {\large Dr. Fernando Villas-Bôas e Prof. Dr. Clóvis
Perin}}
\end{center}  
\vspace{1.5cm}

\begin{center}
{\sffamily Estre trabalho contou com suporte financeiro da FAPESP (processo
2008/09685-3).
 \vspace{1cm}

Campinas-SP

%xxxxx de yyyy
}
\end{center}

\end{titlepage}
\cleardoublepage

% %%%% Capa
% \begin{titlepage}
%   \thispagestyle{empty}
% \begin{center} 
% 
%   
% {\Large {\bfseries \sffamily Luiz Rafael dos Santos}}
% \end{center}
% 
% 
% \vspace{4cm}
% \begin{center}
% \noindent{\LARGE {\sffamily \bfseries  Escolha
% adiada de parâmetros em métodos de pontos \\[2mm] interiores para programação linear}}
% \end{center}
% 
% \vspace{1cm}
% % \setcounter{footnote}{1} \renewcommand{\thefootnote}{\fnsymbol{footnote}}
% \hfill
% \begin{minipage}{9cm}
% {  Tese  apresentada ao Instituto de Matem\'atica, Estat\'{\i}stica
% e Computa\c{c}\~ao Cient\'{\i}fica da Universidade Estadual de Campinas  para a obten\c{c}\~ao do 
% t\'{\i}tulo de {\bfseries DOUTOR em Matem\'atica Aplicada}.} \\
% 
% 
% \end{minipage}
% 
% 
% \vspace{1cm}
% \begin{center}
%  
% {  Orientador: \\ {\large Prof.~Dr.~Aurelio Ribeiro Leite de
% Oliveira}}\\[2mm]
% %Usar se tiver co-orientador
% {  Co-orientadores: \\ {\large Dr.~Fernando da Rocha Villas Bôas e
% Prof.~Dr.~Clóvis Perin}}
% \end{center}  
% 
% 
% \vspace{1cm}
% 
% \begin{center}
% 
% { \large  Universidade Estadual de Campinas \\[1mm] Instituto de Matemática,
% Estatística e Computação Científica \\[1mm] Programa de Pós-graduação em
% Matemática Aplicada}
% 
% \end{center}
% 
% 
% 
% 
% 
% 
% 
% 
% 
% 
% \vfill
% \begin{center}
% { 
% 
% 
% Campinas-SP
% 
% Dezembro de 2012
% }
% \end{center}
% 
% \end{titlepage}
% 
% 


\thispagestyle{empty}
 